\documentclass[12pt]{article}
\usepackage[utf8]{inputenc}
\usepackage[T1]{fontenc}
\usepackage[margin=1in]{geometry}
\usepackage{setspace}
\usepackage{xcolor}
\usepackage[colorlinks=true, urlcolor=blue!60!black, linkcolor=blue!60!black]{hyperref}
\setlength{\parindent}{2em}
\setlength{\parskip}{1.25em}
\renewcommand{\baselinestretch}{1.0}
\begin{document}
\begin{center}
{\LARGE \textit{FDA v. Alliance for Hippocratic Medicine}, 602 U.S. 367 (2024)\par}
\vspace{0.8em}
{\large Case Brief\par}
\end{center}
\vspace{1em}

\section{Memory Jogger}

\textbf{Course Unit:} Standing (Case \#162). \textbf{Core Image:} Pro-life doctors and medical associations who neither prescribe nor use the abortion drug mifepristone cannot establish Article III standing to challenge the FDA's 2016 and 2021 decisions relaxing access to the drug. Federal conscience laws break any causal chain to downstream conscience injuries, there is no Article III doctrine of ``doctor standing'' for attenuated public-health ripple effects, and an advocacy organization ``cannot spend its way into standing'' by diverting resources to oppose a federal policy under \textit{Havens Realty Corp. v. Coleman}. See \href{https://www.supremecourt.gov/opinions/23pdf/602us1r35_bqmd.pdf}{602 U.S. 367 (2024)} (official U.S. Reports preliminary print).

\section{Detailed Case Facts}

Under the Federal Food, Drug, and Cosmetic Act (FDCA), 21 U.S.C. \S~355(d), the U.S. Food and Drug Administration (FDA) reviews new drug applications to ensure that drugs marketed in the United States are safe and effective when used as directed. When additional safeguards are needed, the FDA may impose extra conditions on prescription and distribution under a Risk Evaluation and Mitigation Strategy (REMS), 21 U.S.C. \S~355--1(f)(3), such as prescriber certification, restricted dispensing settings, or patient monitoring (\href{https://www.supremecourt.gov/opinions/23pdf/602us1r35_bqmd.pdf}{602 U.S. at 374--375}).

In 2000, the FDA approved a new drug application for mifepristone tablets, marketed by distributor Danco Laboratories, L.L.C., under the brand name Mifeprex, for use in terminating pregnancies up to seven weeks (\href{https://www.supremecourt.gov/opinions/23pdf/602us1r35_bqmd.pdf}{602 U.S. at 375}). Alongside that approval, the FDA imposed strict conditions of use: only physicians could prescribe or supervise the prescription of Mifeprex, patients had to attend three in-person office visits, and prescribers had to report hospitalizations, blood transfusions, and other serious adverse events to the sponsor, which reported them to the FDA (\href{https://www.supremecourt.gov/opinions/23pdf/602us1r35_bqmd.pdf}{602 U.S. at 375}).

In 2015, Danco Laboratories submitted a supplemental new drug application seeking to update Mifeprex's labeling and ease those restrictions. In 2016, the FDA approved the changes: it extended the approved gestational limit from seven weeks to ten weeks, permitted non-physician healthcare providers such as nurse practitioners to prescribe the drug, reduced the required in-person visits from three to one, and modified prescriber adverse-event reporting to require reporting only of fatalities (\href{https://www.supremecourt.gov/opinions/23pdf/602us1r35_bqmd.pdf}{602 U.S. at 375--376}). In 2019, the FDA approved a generic version of mifepristone under the same conditions of use (\href{https://www.supremecourt.gov/opinions/23pdf/602us1r35_bqmd.pdf}{602 U.S. at 376}). In 2021, drawing on clinical experience during the COVID--19 pandemic, the FDA announced that it would no longer enforce the initial in-person dispensing requirement, effectively allowing telemedicine prescribing and mail-order dispensing (\href{https://www.supremecourt.gov/opinions/23pdf/602us1r35_bqmd.pdf}{602 U.S. at 376}).

Pro-life medical groups had previously filed unsuccessful citizen petitions with the FDA in 2002 (seeking rescission of the 2000 approval) and in 2019 (seeking withdrawal of the 2016 modifications) (\href{https://www.supremecourt.gov/opinions/23pdf/602us1r35_bqmd.pdf}{602 U.S. at 376}). In November 2022---shortly after \textit{Dobbs v. Jackson Women's Health Organization} (2022) and the incorporation of the lead plaintiff, the Alliance for Hippocratic Medicine (an umbrella association whose members are other medical associations rather than individual doctors)---four pro-life medical associations and several individual physicians sued the FDA under the Administrative Procedure Act (APA), challenging the 2000 approval of Mifeprex, the 2019 generic approval, and the 2016 and 2021 regulatory relaxations (\href{https://www.supremecourt.gov/opinions/23pdf/602us1r35_bqmd.pdf}{602 U.S. at 376, 400}; \href{https://harvardlawreview.org/print/vol-138/fda-v-alliance-for-hippocratic-medicine/}{138 Harv. L. Rev. 295, 296 (2024)}). Critically, none of the plaintiff doctors or associations prescribed, manufactured, sold, or used mifepristone, nor did the FDA require them to do or refrain from doing anything (\href{https://www.supremecourt.gov/opinions/23pdf/602us1r35_bqmd.pdf}{602 U.S. at 373--374, 385--386}). Danco Laboratories intervened as a defendant to defend the FDA's regulatory actions (\href{https://www.supremecourt.gov/opinions/23pdf/602us1r35_bqmd.pdf}{602 U.S. at 376--377}).

\section{Procedural History}

\begin{enumerate}
\item \textbf{U.S. District Court for the Northern District of Texas:} In \textit{Alliance for Hippocratic Medicine v. FDA}, 668 F. Supp. 3d 507 (N.D. Tex. 2023), Judge Matthew J. Kacsmaryk held that the plaintiffs possessed Article III standing, concluded they were likely to succeed on the merits of their APA and Comstock Act claims, and issued a preliminary order staying the effective date of the FDA's 2000 approval of mifepristone and all subsequent challenged actions, effectively ordering mifepristone off the market (\href{https://www.supremecourt.gov/opinions/23pdf/602us1r35_bqmd.pdf}{602 U.S. at 377}; \href{https://harvardlawreview.org/print/vol-138/fda-v-alliance-for-hippocratic-medicine/}{138 Harv. L. Rev. at 296--297}).

\item \textbf{Fifth Circuit Partial Stay:} The FDA and Danco appealed and sought an emergency stay. On April 12, 2023, the Fifth Circuit granted the stay in part as to the 2000 approval of Mifeprex (finding challenges to it likely time-barred) but left the injunction in place as to the 2016, 2019, and 2021 actions, which would have reverted Mifeprex to its 2000 restrictions (\href{https://www.supremecourt.gov/opinions/23pdf/602us1r35_bqmd.pdf}{602 U.S. at 377}).

\item \textbf{Supreme Court Emergency Stay:} On April 21, 2023, the Supreme Court stayed the District Court's order in its entirety pending appellate disposition and Supreme Court review, preserving mifepristone's availability under the 2016 and 2021 rules, with Justices Thomas and Alito noting dissent (\href{https://www.supremecourt.gov/opinions/23pdf/602us1r35_bqmd.pdf}{602 U.S. at 377}; \href{https://harvardlawreview.org/print/vol-138/fda-v-alliance-for-hippocratic-medicine/}{138 Harv. L. Rev. at 297}).

\item \textbf{Fifth Circuit Merits Decision:} In \textit{Alliance for Hippocratic Medicine v. FDA}, \href{https://law.justia.com/cases/federal/appellate-courts/ca5/23-10362/23-10362-2023-08-16.html}{78 F.4th 210 (5th Cir. 2023)}, the Fifth Circuit affirmed in part and vacated in part. It held that the individual doctors and medical associations had Article III standing, vacated the District Court's order as to the 2000 and 2019 approvals, and affirmed the preliminary injunction as to the 2016 and 2021 actions as likely arbitrary and capricious under the APA (\href{https://www.supremecourt.gov/opinions/23pdf/602us1r35_bqmd.pdf}{602 U.S. at 377--378}).

\item \textbf{Supreme Court Grant of Certiorari and Decision:} On December 13, 2023, the Supreme Court granted and consolidated the petitions of the FDA (No. 23--235) and Danco Laboratories (No. 23--236) regarding the 2016 and 2021 actions, heard oral argument on March 26, 2024, and unanimously reversed and remanded on June 13, 2024 (\href{https://www.supremecourt.gov/opinions/23pdf/602us1r35_bqmd.pdf}{602 U.S. at 367, 370, 378, 397}; \href{https://supreme.justia.com/cases/federal/us/602/23-235/}{Justia Docket Summary}).
\end{enumerate}

\section{Judicial Votes}

\textbf{Vote:} 9--0 (Unanimous).

\textbf{Opinion of the Court:} Delivered by \textbf{Justice Kavanaugh}, joined in full by Chief Justice Roberts and Justices Thomas, Alito, Sotomayor, Kagan, Gorsuch, Barrett, and Jackson (\href{https://www.supremecourt.gov/opinions/23pdf/602us1r35_bqmd.pdf}{602 U.S. at 370, 372}).

\textbf{Concurring Opinion:} \textbf{Justice Thomas} filed a solo concurring opinion (joining the Court's opinion in full while questioning the constitutionality of associational standing) (\href{https://www.supremecourt.gov/opinions/23pdf/602us1r35_bqmd.pdf}{602 U.S. at 370, 397}).

\textbf{Dissenting Opinions:} None.

\section{Holding}

\textbf{Judgment:} Reversed the judgment of the U.S. Court of Appeals for the Fifth Circuit and remanded for further proceedings consistent with the opinion (\href{https://www.supremecourt.gov/opinions/23pdf/602us1r35_bqmd.pdf}{602 U.S. at 370, 397}).

\textbf{Rule Announced:} Unregulated plaintiffs who do not prescribe, manufacture, sell, or use an FDA-approved drug lack Article III standing to challenge the FDA's relaxation of regulatory restrictions on that drug (\href{https://www.supremecourt.gov/opinions/23pdf/602us1r35_bqmd.pdf}{602 U.S. at 373--374, 385--386}). Specifically: (1) federal conscience protections (such as the Church Amendments, 42 U.S.C. \S~300a--7, which are not overridden by EMTALA) break any causal chain between FDA deregulation and alleged downstream conscience injuries to doctors; (2) alleged downstream monetary, time-diversion, malpractice, or insurance costs from treating third parties' complications are too speculative and attenuated to establish causation, and there is no Article III doctrine of ``doctor standing''; (3) an organization cannot manufacture organizational standing under \textit{Havens Realty Corp. v. Coleman}, 455 U.S. 363 (1982), merely by spending money to study, educate about, or advocate against a government policy that does not directly interfere with the organization's pre-existing core business activities; and (4) the possibility that no one else would have standing does not create standing (\href{https://www.supremecourt.gov/opinions/23pdf/602us1r35_bqmd.pdf}{602 U.S. at 386--396}).

\textbf{Holding vs. Dicta:} The Court's rejection of all three standing theories---conscience injury, downstream economic/``doctor'' standing, and \textit{Havens Realty} resource-diversion organizational standing---as well as its rejection of the ``if not us, who?'' argument and uninjured third-party standing (in footnote 5) are essential to the judgment reversing the Fifth Circuit and constitute binding holdings (\href{https://www.supremecourt.gov/opinions/23pdf/602us1r35_bqmd.pdf}{602 U.S. at 386--396}). By contrast, the Court's discussion of the merits of the FDA's 2016 and 2021 actions was explicitly left unreached; its illustrative hypotheticals regarding EPA emissions rules, speed limits, gun restrictions, firefighters, police officers, and border-state teachers (\href{https://www.supremecourt.gov/opinions/23pdf/602us1r35_bqmd.pdf}{602 U.S. at 391--392}) serve as explanatory reasoning; its observation that ``it is not clear that no one else would have standing'' (\href{https://www.supremecourt.gov/opinions/23pdf/602us1r35_bqmd.pdf}{602 U.S. at 396}) is non-binding dicta; and Justice Thomas's solo critique of associational standing under \textit{Hunt v. Washington State Apple Advertising Commission}, 432 U.S. 333 (1977), is a non-binding concurrence (\href{https://www.supremecourt.gov/opinions/23pdf/602us1r35_bqmd.pdf}{602 U.S. at 397--405}). Notably, however, the Court's endorsement of the Solicitor General's representation that EMTALA does not override individual doctors' federal conscience protections (\href{https://www.supremecourt.gov/opinions/23pdf/602us1r35_bqmd.pdf}{602 U.S. at 388--389}) was an integral step in breaking the causation chain and carries significant doctrinal weight.

\section{Analysis of Opinions}

\subsection{Majority Opinion (Justice Kavanaugh)}

Justice Kavanaugh's opinion for a unanimous Court proceeds in two principal parts: a foundational restatement of Article III standing doctrine (Part II.A) and a systematic rejection of each of the plaintiffs' standing theories (Part II.B) (\href{https://www.law.cornell.edu/supremecourt/text/23-235}{Cornell LII}; \href{https://www.supremecourt.gov/opinions/23pdf/602us1r35_bqmd.pdf}{602 U.S. at 378--397}).

\textbf{Foundational Standing Principles:} Grounding standing in Article III's ``Cases'' and ``Controversies'' text and separation of powers, Kavanaugh invokes Justice Scalia's famous question---``What's it to you?''---to emphasize that federal courts are not an open forum for bystanders with general legal, moral, ideological, or policy objections (\href{https://www.supremecourt.gov/opinions/23pdf/602us1r35_bqmd.pdf}{602 U.S. at 378--382}). A plaintiff must establish (i) a concrete, particularized, and actual or imminent injury in fact; (ii) causation (traceability to the defendant's conduct); and (iii) redressability (\href{https://www.supremecourt.gov/opinions/23pdf/602us1r35_bqmd.pdf}{602 U.S. at 380}). When an unregulated plaintiff challenges the government's regulation (or under-regulation) of someone else, causation is ``ordinarily substantially more difficult to establish'' because it hinges on the choices of independent third parties (\href{https://www.supremecourt.gov/opinions/23pdf/602us1r35_bqmd.pdf}{602 U.S. at 382--383}). The causal chain must be neither too speculative (unpredictable third-party reactions) nor too attenuated (distant ripple effects, even if predictable) (\href{https://www.supremecourt.gov/opinions/23pdf/602us1r35_bqmd.pdf}{602 U.S. at 383}).

\textbf{1. Downstream Conscience Injuries:} The Court agreed that forcing a doctor to provide medical treatment against conscience is a concrete injury in fact (\href{https://www.supremecourt.gov/opinions/23pdf/602us1r35_bqmd.pdf}{602 U.S. at 387}). Even assuming FDA's 2016 and 2021 actions caused more emergency complications, however, federal conscience laws---specifically the Church Amendments, 42 U.S.C. \S~300a--7(c)(1) and (d)---definitively protect doctors from being required to perform or assist in abortions or treatment violating their religious or moral beliefs, even in healthcare deserts (\href{https://www.supremecourt.gov/opinions/23pdf/602us1r35_bqmd.pdf}{602 U.S. at 387--388}). None of the plaintiffs' declarations identified a single instance since 2000 in which a doctor was forced to violate conscience. The Court also agreed with the Solicitor General that EMTALA, 42 U.S.C. \S~1395dd, imposes obligations on hospitals rather than individual physicians and does not override federal conscience protections (\href{https://www.supremecourt.gov/opinions/23pdf/602us1r35_bqmd.pdf}{602 U.S. at 388--389}). In footnote 3, the Court added that mere ``distress at or disagreement with the activities of others'' cannot confer Article III standing (\href{https://www.supremecourt.gov/opinions/23pdf/602us1r35_bqmd.pdf}{602 U.S. at 390 n.3}).

\textbf{2. Downstream Economic Injuries and ``Doctor Standing'':} The doctors' claims of diverted time, increased malpractice risk, and potentially higher insurance premiums lacked record support and were both speculative and impermissibly attenuated (\href{https://www.supremecourt.gov/opinions/23pdf/602us1r35_bqmd.pdf}{602 U.S. at 390--391}). Because virtually all FDA-approved drugs carry side effects, recognizing ``doctor standing'' would allow physicians to challenge any drug approval or general safety deregulation---from EPA emissions rollbacks to higher speed limits or repealed gun restrictions---and could not be principledly cabined to doctors, inviting suits by firefighters, police officers, and teachers (\href{https://www.supremecourt.gov/opinions/23pdf/602us1r35_bqmd.pdf}{602 U.S. at 391--392}). In footnote 5, the Court also rejected third-party standing on behalf of patients because third-party standing still requires the litigant itself to suffer an injury in fact (\href{https://www.supremecourt.gov/opinions/23pdf/602us1r35_bqmd.pdf}{602 U.S. at 393 n.5}).

\textbf{3. Organizational Standing and Narrowing \textit{Havens Realty}:} The medical associations argued that FDA's actions forced them to spend resources conducting studies, filing citizen petitions, and educating the public. Kavanaugh flatly rejected this resource-diversion theory: an organization ``cannot spend its way into standing'' (\href{https://www.supremecourt.gov/opinions/23pdf/602us1r35_bqmd.pdf}{602 U.S. at 394}). The Court sharply cabined \textit{Havens Realty Corp. v. Coleman}, 455 U.S. 363 (1982), explaining that Havens Realty's racial steering directly gave false information to the plaintiff's employees and ``directly affected and interfered with HOME's core business activities---not dissimilar to a retailer who sues a manufacturer for selling defective goods to the retailer'' (\href{https://www.supremecourt.gov/opinions/23pdf/602us1r35_bqmd.pdf}{602 U.S. at 395}). Describing \textit{Havens} as ``an unusual case'' that the Court has been careful not to extend beyond its context, the Court concluded that FDA's deregulation imposed no comparable impediment to the associations' advocacy operations (\href{https://www.supremecourt.gov/opinions/23pdf/602us1r35_bqmd.pdf}{602 U.S. at 395--396}). Finally, the Court reiterated that even if no one else would have standing, Article III leaves some contested policy questions to the political, legislative, and regulatory branches (\href{https://www.supremecourt.gov/opinions/23pdf/602us1r35_bqmd.pdf}{602 U.S. at 396--397}).

\subsection{Concurrence (Justice Thomas)}

Justice Thomas joined the Court's opinion in full but wrote separately to urge the Court, in an appropriate future case, to reconsider both third-party standing and \textbf{associational standing} under \textit{Hunt v. Washington State Apple Advertising Commission}, 432 U.S. 333 (1977) (\href{https://www.supremecourt.gov/opinions/23pdf/602us1r35_bqmd.pdf}{602 U.S. at 397--399}). Relying on historical common-law limits on the judicial power, Thomas argued that associational standing violates Article III in two ways: (1) it relaxes the injury requirement by allowing an uninjured organization to sue vicariously based on a single member's injury (pointing out that the Alliance is an association of associations, two degrees removed from any doctor, and citing amicus Professor F. Andrew Hessick's example of AARP's nearly 38 million members); and (2) it distorts redressability and traditional equity limits by driving courts to issue ``universal injunctions'' benefitting nonparties (\href{https://www.supremecourt.gov/opinions/23pdf/602us1r35_bqmd.pdf}{602 U.S. at 399--402}). He added that associational standing subverts Federal Rule of Civil Procedure 23's class-action safeguards, creates asymmetrical preclusion (potentially giving members ``two bites at the apple''), and arose ``seemingly by accident'' from a one-paragraph rehearing denial in \textit{National Motor Freight Traffic Assn., Inc. v. United States}, 372 U.S. 246 (1963) (\href{https://www.supremecourt.gov/opinions/23pdf/602us1r35_bqmd.pdf}{602 U.S. at 402--405}).

\subsection{Dissents}

There were no dissenting opinions.

\section{Examples: Future Applications}

\begin{enumerate}
\item \textbf{Same-Side (No Standing): Voting-Rights Nonprofits Challenging State Voter-Roll Rules.} Suppose civic advocacy groups sue to enjoin a state law modifying voter-registration cancellation procedures, alleging that the law frustrates their mission and forces them to divert staff and funding to retrain volunteers and educate voters. Under \textit{Alliance for Hippocratic Medicine}, these organizations lack Article III standing because they cannot ``spend [their] way into standing'' by voluntarily reallocating resources to educate or advocate against a policy that does not directly obstruct their pre-existing core operations. Indeed, the Ninth Circuit applied \textit{Alliance for Hippocratic Medicine} to reach this exact result in \href{https://www.govinfo.gov/content/pkg/USCOURTS-ca9-22-16490/pdf/USCOURTS-ca9-22-16490-0.pdf}{\textit{Arizona Alliance for Retired Americans v. Mayes}} (9th Cir. 2024; en banc 2026).

\item \textbf{Same-Side (No Standing): Emergency Physicians Challenging Deregulation of Motorized Scooters.} Suppose a coalition of emergency room doctors sues the Consumer Product Safety Commission after it repeals helmet and speed-limiter standards for electric scooters, alleging that more injured riders will crowd their emergency departments, divert time from other patients, and raise malpractice exposure. Under \textit{Alliance for Hippocratic Medicine}, the doctors lack standing: they are unregulated parties challenging the under-regulation of third parties, and there is no Article III doctrine of ``doctor standing'' for attenuated public-safety ripple effects (\href{https://www.supremecourt.gov/opinions/23pdf/602us1r35_bqmd.pdf}{602 U.S. at 391--392}).

\item \textbf{Opposite-Side (Standing Exists): Competing Pharmaceutical Manufacturer Challenging Unlawful Approval.} Suppose a pharmaceutical company that markets an FDA-approved rival medication for the same medical indication sues the FDA under the APA, alleging that the agency unlawfully waived statutory safety-testing requirements to approve a cheaper competing drug. Unlike the pro-life doctors in \textit{Alliance for Hippocratic Medicine}, the rival manufacturer suffers a direct, predictable downstream economic injury as a market competitor---a classic category of unregulated-party standing that Justice Kavanaugh expressly preserved (\href{https://www.supremecourt.gov/opinions/23pdf/602us1r35_bqmd.pdf}{602 U.S. at 384--385}).

\item \textbf{Opposite-Side (Standing Exists): Doctor Directly Regulated by a Hospital Safety Mandate or Stripped of Conscience Protections.} Suppose a state enacts a statute revoking state conscience exemptions and threatening medical-license revocation against any obstetrician who refuses to perform a specified procedure, or suppose the FDA imposes a mandatory reporting and training rule directly on emergency room physicians. In either case, the doctor is directly regulated by the government or faces a credible threat of being forced to act against conscience, satisfying both injury in fact and direct causation (\href{https://www.supremecourt.gov/opinions/23pdf/602us1r35_bqmd.pdf}{602 U.S. at 382, 387, 392 n.4}).

\item \textbf{Fence-Sitter (Close Case): Direct-Service Fair-Housing or Immigration Legal-Aid Nonprofit Facing Direct Operational Obstruction.} Suppose a nonprofit that both advocates for immigrants and provides direct legal representation challenges a federal agency policy that bars its attorneys from accessing detained clients or systematically provides false case-status information to its staff. Unlike the medical associations in \textit{Alliance for Hippocratic Medicine}, the group provides direct services rather than pure issue advocacy and can analogize to \textit{Havens Realty}'s ``defective goods'' impairment of core business activities (\href{https://www.supremecourt.gov/opinions/23pdf/602us1r35_bqmd.pdf}{602 U.S. at 395}). Yet lower courts after \textit{Alliance for Hippocratic Medicine}---such as the en banc Fifth Circuit in \href{https://law.justia.com/cases/federal/appellate-courts/ca5/24-50149/24-50149-2026-04-24.html}{\textit{United States v. Texas}} (5th Cir. 2026)---have scrutinized whether a policy merely makes a nonprofit's work costlier or actually obstructs its pre-existing core operations, leaving direct-service nonprofits in a sharply contested doctrinal gray zone.
\end{enumerate}

\section{Critique}

\textbf{Originalist and Textualist Perspective:} From an originalist and textualist standpoint, the majority opinion is widely praised as a disciplined, clear restatement of Article III's separation-of-powers limits that keeps federal courts from becoming roving policy monitors (\href{https://reason.com/volokh/2024/06/13/the-standing-analysis-in-fda-v-alliance-for-hippocratic-medicine/}{Josh Blackman, Volokh Conspiracy (June 13, 2024)}). Conservative commentators welcomed the Court's sharp narrowing of \textit{Havens Realty} ``diversion of resources'' standing and its footnote 3 rejection of ``distress'' or ``offended observer'' standing, arguing both doctrines strayed from Article III's text and history (\href{https://reason.com/volokh/2024/06/13/the-standing-analysis-in-fda-v-alliance-for-hippocratic-medicine/}{Volokh Conspiracy}). Justice Thomas's concurrence pushed the originalist critique even further, contending that the majority stopped short by leaving \textit{Hunt}'s associational-standing framework intact despite its lack of common-law pedigree and its connection to universal injunctions (\href{https://www.supremecourt.gov/opinions/23pdf/602us1r35_bqmd.pdf}{602 U.S. at 399--405}). At the same time, some conservative observers noted that the Fifth Circuit's conscience-standing analysis had relied on the federal government's earlier, stricter EMTALA litigation posture in \textit{Texas v. Becerra}, and that Solicitor General Prelogar's broad oral-argument concessions regarding federal conscience protections effectively mooted that theory while cementing robust conscience rights for future cases (\href{https://reason.com/volokh/2024/06/13/the-standing-analysis-in-fda-v-alliance-for-hippocratic-medicine/}{Volokh Conspiracy}).

\textbf{Progressive and Doctrinal Perspective:} Progressive and doctrinal critics view the unanimous outcome with caution rather than celebration. As the \href{https://harvardlawreview.org/print/vol-138/fda-v-alliance-for-hippocratic-medicine/}{Harvard Law Review (138 Harv. L. Rev. at 295, 298--304)} observed, the decision resolved only a threshold procedural barrier and left the substantive administrative-law and Comstock Act challenges to mifepristone open to future suits by states such as Missouri, Kansas, and Idaho---particularly in a post-\textit{Loper Bright} and post-\textit{Corner Post} landscape where agency deference has been eliminated and APA limitations periods lengthened. Critics also charge the Court with doctrinal inconsistency and strategic malleability in standing law, contrasting its rigorous causation demands here with its willingness to find standing on tenuous theories in \textit{Biden v. Nebraska} (2023) and \textit{303 Creative LLC v. Elenis} (2023) (\href{https://harvardlawreview.org/print/vol-138/fda-v-alliance-for-hippocratic-medicine/}{138 Harv. L. Rev. at 299--300}). Moreover, both progressive and conservative commentators warn that the unanimous narrowing of \textit{Havens Realty} and Justice Thomas's warning shot against associational standing will substantially curtail civil-rights, voting-rights, environmental, and fair-housing nonprofits from challenging unlawful government policies in federal court (\href{https://reason.com/volokh/2024/06/13/the-standing-analysis-in-fda-v-alliance-for-hippocratic-medicine/}{Volokh Conspiracy}; \href{https://www.scotusblog.com/2024/06/supreme-court-preserves-access-to-abortion-pill/}{SCOTUSblog}).

\textbf{Underlying Competing Values:} Beneath the technical causation analysis lies a deep structural tension between \textbf{judicial restraint and democratic accountability} on one hand, and \textbf{meaningful judicial review of administrative legality} on the other. Strict standing rules prevent ideological litigants from using nationwide injunctions to hijack executive-branch scientific judgments, channeling policy disputes back to Congress, the President, and the electorate (\href{https://www.supremecourt.gov/opinions/23pdf/602us1r35_bqmd.pdf}{602 U.S. at 380, 396--397}). Yet when an agency deregulates---making a product easier rather than harder to obtain---strict traceability rules can insulate potentially unlawful deregulation from any judicial check at all, forcing the legal system to confront how much unreviewable executive discretion Article III's separation of powers is willing to tolerate.

\section{Key Quotations}

\begin{quote}
``The injury in fact requirement prevents the federal courts from becoming a `vehicle for the vindication of the value interests of concerned bystanders.' \textit{Allen}, 468 U. S., at 756 (quotation marks omitted). An Article III court is not a legislative assembly, a town square, or a faculty lounge. Article III does not contemplate a system where 330 million citizens can come to federal court whenever they believe that the government is acting contrary to the Constitution or other federal law.''
\end{quote}
\noindent --- Justice Kavanaugh, Opinion of the Court, \href{https://www.supremecourt.gov/opinions/23pdf/602us1r35_bqmd.pdf}{\textit{FDA v. Alliance for Hippocratic Medicine}, 602 U.S. 367, 382 (2024)} (also available via \href{https://www.law.cornell.edu/supremecourt/text/23-235}{Cornell LII}).

\begin{quote}
``Federal law fully protects doctors against being required to provide abortions or other medical treatment against their consciences---and therefore breaks any chain of causation between FDA's relaxed regulation of mifepristone and any asserted conscience injuries to the doctors.''
\end{quote}
\noindent --- Justice Kavanaugh, Opinion of the Court, \href{https://www.supremecourt.gov/opinions/23pdf/602us1r35_bqmd.pdf}{\textit{FDA v. Alliance for Hippocratic Medicine}, 602 U.S. 367, 390 (2024)} (also available via \href{https://supreme.justia.com/cases/federal/us/602/23-235/}{Justia}).

\begin{quote}
``Stated otherwise, there is no Article III doctrine of `doctor standing' that allows doctors to challenge general government safety regulations. Nor will this Court now create such a novel standing doctrine out of whole cloth.''
\end{quote}
\noindent --- Justice Kavanaugh, Opinion of the Court, \href{https://www.supremecourt.gov/opinions/23pdf/602us1r35_bqmd.pdf}{\textit{FDA v. Alliance for Hippocratic Medicine}, 602 U.S. 367, 391 (2024)} (also available via \href{https://www.law.cornell.edu/supremecourt/text/23-235}{Cornell LII}).

\begin{quote}
``But an organization that has not suffered a concrete injury caused by a defendant's action cannot spend its way into standing simply by expending money to gather information and advocate against the defendant's action. An organization cannot manufacture its own standing in that way.''
\end{quote}
\noindent --- Justice Kavanaugh, Opinion of the Court, \href{https://www.supremecourt.gov/opinions/23pdf/602us1r35_bqmd.pdf}{\textit{FDA v. Alliance for Hippocratic Medicine}, 602 U.S. 367, 394 (2024)} (also available via \href{https://www.supremecourt.gov/opinions/23pdf/23-235_n7ip.pdf}{Supreme Court Slip Opinion}).

\begin{quote}
``Associational standing seems to run roughshod over this traditional understanding of the judicial power. Our doctrine permits an association to have standing based purely upon a member's injury, not its own. If a single member of an association has suffered an injury, our doctrine permits that association to seek relief for its entire membership---even if the association has tens of millions of other, non-injured members.''
\end{quote}
\noindent --- Justice Thomas, concurring, \href{https://www.supremecourt.gov/opinions/23pdf/602us1r35_bqmd.pdf}{\textit{FDA v. Alliance for Hippocratic Medicine}, 602 U.S. 367, 399--400 (2024)} (also available via \href{https://www.law.cornell.edu/supremecourt/text/23-235}{Cornell LII}).

\end{document}
